\begin{solapartat}
\begin{minipage}[t]{0.42\linewidth}
\punts{0,3} Si l'ió no es desvia de la trajectòria rectilínia, s'ha de complir que la resultant de forces sigui zero sobre l'ió He, la \textbf{força elèctrica} ha de tenir el mateix mòdul que la força magnètica i sentit oposat.
\end{minipage}\hfill
\begin{minipage}[t]{0.52\linewidth}
\vspace{0pt}
\figura[0.75]{sol_fig1.png}
\end{minipage}

\medskip
\punts{0,1} I atès que l'ió és positiu, \textbf{el camp elèctric} tindrà el mateix sentit que la força elèctrica.

\punts{0,2} $\vec{E} = -2,00\cdot 10^{5}\,\vec{\jmath}\ N/C$

\punts{0,2} $\begin{aligned}[t] F_m &= F_e\\ qvB &= qE \end{aligned}$

\punts{0,2} $B = \dfrac{E}{v} = \dfrac{2,00\times 10^{5}}{3,2\times 10^{5}} = 0,625\ T$
\end{solapartat}

\begin{solapartat}
\punts{0,4} $F_m = ma_c$

\punts{0,2} $qvB = m\,\dfrac{v^2}{R} \rightarrow R = \dfrac{mv^2}{qvB} = \dfrac{mv}{qB}$

\punts{0,4} $R = \dfrac{6,68\times 10^{-27}\cdot 3,2\times 10^{5}}{1,6\cdot 10^{-19}\cdot 0,625} = 0,0214\ m$
\end{solapartat}
